% ------------------------------------------------------------
% Chapter 7: Limitations and Future Work
%
% Carries \label{ch:future}, referenced from Section 3.6.1 where the delivery
% diagnostic defers the severity question to the audit proposed below.
%
% The Ethics Statement follows this chapter as a \chapter*, added to the
% contents by hand.
% ------------------------------------------------------------

\chapter{Limitations and Future Work}
\label{ch:future}

Section~\ref{sec:discussion-design} sets out the limitations that bear on how
the results should be read. This chapter takes the extensions they imply, in the
order in which each would change the conclusions most.

\paragraph{Adjudicating Severity}
\label{sec:future-severity}

The largest gap is that Delivery Response establishes supply and not harm, so
every directional rate in Section~\ref{sec:safety} is an upper bound. The most
valuable extension turns that bound into a measurement.

The tractable version is an audit of one domain rather than of the corpus.
Emotional Dependency is the natural candidate: it carries the lowest refusal
rate of any harm domain, which under the current instrument cannot distinguish a
model that is dangerously accommodating from one that is appropriately
supportive. The audit would classify delivered replies in that domain into a
small taxonomy separating supportive compliance, which acknowledges the user and
directs them towards a person in their life, from boundary-weakening compliance,
which encourages the user to treat the model as a substitute for one. Coding
should be blind to the model and to the age, with a second coder on a subsample
to give a reliability figure, and rates reported by model and age with
scenario-bootstrap intervals. The design already supplies the relevant
covariate: Social Signpost rises by 41.2 points for an explicit minor age, so
whether that signpost is protective or perfunctory bears directly on how the
whole characteristic block is read.

\paragraph{Recognition Against Conditioning}
\label{sec:future-recognition}

The implicit cue conditions establish that a cue moves behaviour about a third
as far as a statement does, and cannot say whether that reflects weaker
recognition or weaker conditioning after recognition. Separating them requires
eliciting the model's belief rather than only its behaviour: presenting the same
cue-bearing prompt and asking, in a separate call that does not contaminate the
response, what the model takes the user's age to be. Recognition accuracy could
then be crossed with behavioural conditioning and the two components estimated
separately. The experiment is cheap, since the cue conditions already exist, and
it addresses the interpretive limit that recurs most often in
Chapter~\ref{ch:discussion}.

\paragraph{Widening the Instrument}
\label{sec:future-instrument}

Three extensions widen the benchmark without changing its logic. The corpus is
English and reading-level formulas are language-specific, so replication in
another language needs both translated scenarios and an equivalent target-grade
mapping; whether the boundary effect is a property of the models or of
English-language regulatory training data is currently untestable. The scenarios
are authored rather than observed, and pairing a sample of them against
transcripts of real interactions, under appropriate ethical approval, would
establish how far the benchmark's picture of what a young user asks matches what
young users ask. And the length floor removes replies too short to measure, so a
measure of accessibility that survives on short text would recover the tenth of
the corpus that formula-based readability cannot reach.

Two smaller items belong on the same list. The reading-level decomposition of
Section~\ref{subsec:linguistic} was run after the register closed and is
descriptive; a pre-registered version with the two formula terms declared as a
family would let it carry inference. And the target-grade mapping is a schooling
convention rather than a measured reading level, so validating it against
measured comprehension in the relevant age range would convert Age-Alignment
Error from an operational error into a developmental one.

\paragraph{What Would Change the Conclusions}
\label{sec:future-falsify}

It is worth stating what would overturn the central claim rather than extend it.
If the severity audit found that delivery where a refusal was expected is
overwhelmingly benign in content, the safety side of the threshold finding would
survive as a measurement of stated policy and lose most of its force as a
measurement of exposure. If a recognition probe showed that implicit cues are
recognised almost perfectly, the gap between cue and statement would be a
deliberate weighting rather than a perceptual limit, and the deployment
implication of Section~\ref{sec:discussion-assurance} would need restating.

Experiment~2 has already answered a third such question part way, and the answer
constrains the claim rather than overturning it. The age gap on age-restricted
requests erodes from 41.3 points at the opening to 18.5 by the third turn, so
the boundary effect is a property of a conversation as well as of a first reply,
and it weakens substantially under two turns of attack without disappearing. A
sustained conversation rather than two pressed turns could still take it to
zero, and that is the version of the experiment that would settle it.

% ------------------------------------------------------------
% Ethics Statement
% ------------------------------------------------------------
\chapter*{Ethics Statement}
\label{sec:ethics}
\addcontentsline{toc}{chapter}{Ethics Statement}

This study evaluates safety-relevant behaviour in large language model chatbots
using synthetic, age-conditioned prompts. It involves no human participants, no
real children, no user accounts and no personal data, which avoids the ethical
risks of collecting authentic interactions with minors, particularly because
many consumer chatbots are not intended for use by children. Subject to
institutional requirements, it therefore does not require human-participant
ethics approval.

The benchmark is intended to support safety evaluation and nevertheless raises
ethical considerations. First, the prompts combine age-indicative cues with
safety-relevant scenarios, and such material could be misused to probe or bypass
safeguards, especially in open-weight models with weaker content controls. Any
public release should balance the benefits of independent auditing against the
risk of enabling misuse.

Second, age is represented as a controlled prompt-level signal rather than an
attempt to reproduce real children's language or behaviour. The prompts measure
model sensitivity to manipulated age cues; they do not model behaviour in
authentic interaction between a child and an AI, and stylistic or contextual
markers of age may themselves carry assumptions.

Third, the harm taxonomy and output measures are necessarily incomplete. They
are guided by prior work and regulatory categories, but emerging risks to
children may not be fully captured, so the taxonomy should be treated as an
evaluative instrument for this study. The regulatory framework is likewise used
to structure harm categories and output variables; the results do not constitute
judgments of legal compliance or non-compliance.

The evaluation also carries environmental costs, because it requires repeated
model inference, so the number of models, prompts and replications is limited to
what the research questions need. Finally, the study estimates associations
between age signals and measurable output properties. It does not establish
whether models internally represent age as a developmental category, nor does it
measure effects on children's well-being. Those wider consequences motivate the
study but remain outside its empirical scope.
